\documentclass[12pt]{article}
\pdfoutput=1

\usepackage[a4paper,margin=24mm]{geometry}
\usepackage{amsmath,amssymb,bm}
\usepackage{booktabs}
\usepackage{graphicx}
\usepackage{microtype}
\usepackage[numbers,sort&compress]{natbib}
\usepackage[hidelinks]{hyperref}
\usepackage{xcolor}

\def\restoreVfifteendigitcatcodes{\catcode`1=12\catcode`5=12}
\catcode`5=11
\catcode`1=11
% Generated from V15 aggregate results; do not edit manually.
\newcommand{\V15Trajectories}{48}
\newcommand{\V15Decisions}{34,560}
\newcommand{\V15ClustersPerModel}{6}
\newcommand{\V15GPUHours}{49.41}
\newcommand{\V15OriginalGPUHours}{19.19}
\newcommand{\V15ReconstructionGPUHours}{49.41}
\newcommand{\V15OrphanFormalDecisions}{197}
\newcommand{\V15OrphanFormalGPUHours}{0.328}
\newcommand{\V15UnrecordedInfrastructureCalls}{1}
\newcommand{\V15HOne}{42.263}
\newcommand{\V15HOneCI}{24.316 to 59.303}
\newcommand{\V15HTwo}{0.05438}
\newcommand{\V15HTwoCI}{0.03459 to 0.07548}
\newcommand{\V15HThree}{0.04845}
\newcommand{\V15HThreeCI}{0.02190 to 0.07526}
\newcommand{\V15HFour}{52.541}
\newcommand{\V15HFourCI}{35.406 to 68.673}
\newcommand{\V15HOneDisposition}{supported}
\newcommand{\V15HTwoDisposition}{supported}
\newcommand{\V15HThreeDisposition}{supported}
\newcommand{\V15HFourDisposition}{supported}
\newcommand{\V15QwenCorrQuench}{0.00238}
\newcommand{\V15QwenCorrQuenchCI}{-0.01750 to 0.02277}
\newcommand{\V15GraniteCorrQuench}{-0.02032}
\newcommand{\V15GraniteCorrQuenchCI}{-0.05486 to 0.00947}
\newcommand{\V15QwenTauMemory}{-2.46995}
\newcommand{\V15QwenTauMemoryCI}{-5.79214 to 0.17724}
\newcommand{\V15GraniteTauMemory}{-4.14109}
\newcommand{\V15GraniteTauMemoryCI}{-19.88286 to 11.60067}
\newcommand{\V15QwenBinderQuench}{-4.80504}
\newcommand{\V15QwenBinderQuenchCI}{-7.42824 to -2.20832}
\newcommand{\V15GraniteBinderQuench}{0.16643}
\newcommand{\V15GraniteBinderQuenchCI}{-0.09447 to 0.52063}

\restoreVfifteendigitcatcodes
\newcommand{\E}{\mathbb E}
\newcommand{\KL}{D_{\mathrm{KL}}}
\newcommand{\figroot}{figures}
\newcommand{\resultfigure}[4]{%
\begin{figure}[htbp]\centering
\IfFileExists{\figroot/#1}{\includegraphics[width=#2]{\figroot/#1}}{\fbox{\parbox[c][45mm][c]{0.88\textwidth}{\centering Figure generated after formal analysis:\\[2mm]\texttt{\detokenize{#1}}}}}
\caption{#3}\label{#4}\end{figure}}



\title{Statistical mechanical characterization of memory and quench response in state-separated LLM-agent networks}

\author{Alexander V. Mantzaris\\
  School of Data, Mathematical, and Statistics Sciences \\
  University of Central Florida\\
  Orlando, FL 32816, USA\\
  \texttt{alexander.mantzaris@ucf.edu}}

\date{}


\begin{document}
\maketitle

\begin{abstract}
Populations of interacting large-language-model (LLM) agents can develop
collective dynamics that are not visible in individual responses or final
averages. Statistical mechanics provides a principled framework for measuring
these dynamics through order, fluctuations, correlations, response,
persistence, and temporal asymmetry. We treat networks of locally informed LLM
agents as finite interacting stochastic systems with private state, bounded
memory, and networked communication.

We study modular networks of 16 Qwen and Granite agents under
random sequential updates. The study combines a corrected analysis of an
earlier experiment, without altering its trajectories, with a prospectively
specified cross-model experiment comprising 48 trajectories, 34,560 local
decisions, and six independent graph environment clusters per model. We
reverse and restore an external field while comparing perturbed systems with
matched nominal trajectories. To isolate memory, we compare persistent
histories with current state Markovized prompts and format-matched scrambled
histories.

All four prespecified tests met their inferential criteria. Field reversal
produced a cross model response of 42.3 within model distance units (95\% CI
$[24.3,59.3]$), and the fixed restoration contrast was 52.5 units (95\% CI
$[35.4,68.7]$). Persistent memory increased bias adjusted and path reversal
divergence by 0.054 nats per attempted update relative to Markovized histories
(95\% CI $[0.035,0.075]$) and by 0.048 relative to scrambled histories (95\%
CI $[0.022,0.075]$). The latter comparison indicates that temporally
meaningful history matters beyond the presence of additional prompt content,
with especially pronounced and consistent effects for Granite. An
out of sample kinetic surrogate identified which aspects of the response admit
a simple local closure and which remain model and history dependent.

These results show that statistical mechanical observables recover collective
response, memory effects, restoration dynamics, and reduced structure that
population means discard. The interpretation remains effective and
finite size  decoding temperature is not physical temperature, reference
energy is not literal energy, and projected path reversal divergence is not
exact thermodynamic entropy production.
\end{abstract}


\noindent\textbf{Keywords:} statistical mechanics; agents; language-model
agents; quench dynamics; entropy rate; multi-information; temporal asymmetry.

\section{Introduction}
\label{sec:introduction}

Large language models are increasingly studied not only as isolated
prompt response systems but as components of interacting agent populations
\cite{wang2024agents,guo2024multiagents}.  Once agents exchange messages,
retain histories, and act through a network, an output from one local decision
can become part of another agent's next input.  This feedback permits
amplification, cascades, coordination, fragmentation, lock-in, delayed
response, and incomplete recovery even when every component uses fixed model
weights.  The capability of one model instance therefore does not determine
the trajectory of the group, because repeated circulation can create modes of
organization that no single generation exhibits.  The scientific question
becomes what dynamical organization emerges when many stochastic
language-model instances interact
over time, a question that matters for the analysis, design, reliability, and
control of multi-agent AI systems.

At the component level, an LLM agent places a generative model inside a
recurrent loop of observation, decision, action, and environmental feedback.
Its future inputs can contain its own earlier outputs, received messages, and
changes caused by earlier actions.  Population trajectories can consequently
depend on message content, network position, update order, retained context,
and decoding randomness.  Conventional benchmarks can characterize useful
local capabilities, but they do not reveal whether network feedback produces
stable organization, amplifies a perturbation, stores hidden history, or
returns the population to its earlier regime.  Those are properties of the
interacting process rather than of any response considered alone.

Recent agent architectures establish both the feasibility of such processes
and the source of their complexity.  Role-playing agents cooperate through
natural-language exchange \cite{li2023camel}, while generative agents combine
observation, planning, reflection, and persistent natural-language memory in a
simulated social environment \cite{park2023generative}.  These systems show
that retrieved history and communication can materially condition later
actions.  Yet architecture and application studies are commonly evaluated by
task completion or behavioral plausibility.  Those objectives do not by
themselves specify a controlled stochastic process, distinguish private state
from observed state, or identify which population coordinates retain
dynamical information.

Computational social simulation and population experiments make the same
measurement problem visible from another direction.  Conditioned language
models reproduce portions of survey-response distributions
\cite{argyle2023}, and simulated samples replicate some established behavioral
experiments while exhibiting model specific distortions, in \cite{aher2023} the
broader opportunities and validity limits of LLM-empowered agent-based
simulation are reviewed in \cite{gao2024abm}.  Networked agents have generated
consensus and fragmentation in opinion experiments \cite{chuang2024}, shared
naming conventions and collective bias \cite{flintashery2025}, and
model-dependent cooperation and coordination in repeated games
\cite{akata2025}.  Together, these studies establish that interaction can
produce reproducible group level outcomes, while providing lessons against treating
semantic possibility as a dynamical explanation or an LLM population as a
faithful proxy for human behavior.

When agents interact through a network, their collective behavior depends not
only on how individual agents update, but also on which agents can communicate
with one another. The interaction graph defines the routes along which
information, influence, and external perturbations can spread. A densely
connected network may transmit a local change throughout the population,
whereas a modular network may initially confine that change to a single
community. Statistical mechanics and network science provide tools for
relating these microscopic interaction patterns to macroscopic collective
behavior \cite{albert2002,boccaletti2006}. Previous studies of cooperative
processes on networks have also shown that network topology and finite
population size can substantially affect the collective behavior that is
observed \cite{dorogovtsev2008}.
Population averages alone cannot reveal this internal organization. Two agent
populations can have the same mean level of order while differing in their
fluctuations, correlations, community structure, and response to perturbation.
We therefore measure how strongly agent states are correlated as a function of
their distance on the interaction graph. These graph-distance correlations
indicate whether collective organization remains localized among neighboring
agents or extends across the wider network. Such network-aware observables
complement population means and provide a more complete description of how
interacting agent groups organize and respond.

More broadly, statistical mechanics supplies a disciplined framework for
connecting microscopic interactions to macroscopic behavior. Agent-based
models ask how local rules and heterogeneous interactions generate
system-level outcomes \cite{bonabeau2002}, while statistical physics has
developed quantitative descriptions of opinion, cultural, language, and
consensus dynamics \cite{castellano2009,baronchelli2018}. In a conventional
stochastic model, the microscopic rule may be simple: a spin flips with a
specified probability, or an agent adopts the state of a neighbor. Such local
rules can nevertheless generate collective relaxation \cite{glauber1963},
stationary behavior constrained by a potential \cite{blume1993}, or distinct
forms of organization when interactions are nonreciprocal
\cite{fruchart2021}.

An LLM agent makes the same micro-to-macro question more challenging because
its local update is not described by a simple flip probability. When selected
for an update, an agent interprets a language prompt containing its current
state, messages from neighboring agents, an external field, and, when
available, a private interaction history. For example, an agent receiving
several neighboring messages in favor of one action may follow that local
majority, resist it because of its remembered history, or produce a different
choice under another stochastic generation. The microscopic transition is
therefore language-conditioned, history-dependent, and high-dimensional, even
when the recorded belief and action are categorical. Statistical-mechanical
observables make the resulting population dynamics tractable without requiring
a complete analytic model of every language-mediated decision. They allow us
to ask whether local updates create population order, whether perturbations
produce a coordinated response, whether correlations extend across the
network, and whether private memory changes the apparent temporal dynamics.

For a finite population observed over finite time windows, fluctuations,
correlations, response distributions, and variation across independent
realizations contain information that a mean trajectory discards. Statistical
mechanics can therefore be used here as a measurement framework for collective
LLM-agent behavior without assuming literal particles, a thermodynamic limit,
or the existence of a phase transition.


Controlled perturbation is important because nominal or steady observations
underdetermine dynamics.  Similar baseline averages could describe a robustly
organized system, a fragile one, a system that responds strongly and
recovers, or one with delayed or incomplete recovery.  Reversing an external
field and then restoring it probes finite-horizon sensitivity, response,
relaxation, and recovery directly.  Comparing that trajectory with a matched
nominal arm also separates intervention-associated change from ordinary
evolution over the same observation window.  A quench is used here as a
dynamical intervention, not as evidence that the network is an equilibrium
material.

Persistent history creates a distinct coarse-graining problem.  The complete
agent process may include private textual state that is absent from a measured
belief then action projection, so the observed population trajectory need not be
Markov even when the augmented simulator state is.  Comparing persistent and
Markovized prompts tests the effect of retaining history; comparing genuine
history with a scrambled, format-matched control helps separate temporally
related information from the mere presence of additional prompt material.
Temporal correlation then measures persistence, while path reversal tests
ordering in the projected trajectory.  State-space currents and trajectory
reversal have exact thermodynamic interpretations in appropriate fully
observed Markov systems \cite{seifert2012,roldan2012}, but coarse graining can
hide dissipation and make inferred bounds estimator dependent
\cite{kawaguchi2013,busiello2019}.  The path statistic used here is therefore
projected temporal asymmetry, not literal entropy production.

Recent LLM-population studies  approach these questions with
statistical-physics concepts, group size can alter collective misalignment
\cite{flint2025}, topology and self social weighting can change conformity
\cite{han2026}, intrinsic fields can outweigh fitted neighbor coupling
\cite{denobili2026}, debate noise can produce model dependent finite size
crossovers \cite{okawa2026}, and high dimensional interaction data can admit
compact dynamical regimes \cite{el2026}.  This progress narrows the gap, but
the joint problem remains insufficiently resolved that controlled state and
information boundaries, field reversal and restoration, genuine and placebo
history interventions, spatial and temporal diagnostics, cluster-level
inference, and an out-of-sample closure test are rarely examined together.
A task score or single consensus coordinate cannot answer all of these
questions.


No single measurement is sufficient to describe how a population of LLM agents
organizes. For example, two populations can have the same average belief while
differing substantially in their internal behavior. One may be relatively
uniform, while the other may contain distinct communities or fluctuate between
different collective states. We therefore use several complementary
measurements. Belief and action order quantify the extent of population
alignment, while their overlap indicates whether the agents' stated beliefs
are reflected in their committed actions. Fluctuations, entropy, dependence,
and effective compatibility then show whether this alignment is stable,
variable, coordinated, or structured by interactions among neighboring
agents. Graph-distance correlations reveal whether changes remain local or
extend across the network. Autocorrelation and Binder statistics provide
additional finite-window descriptions of persistence and of the distribution
of collective order. These diagnostics characterize the observed trajectories
but are not treated as evidence of equilibrium correlation times, critical
slowing down, or a phase transition.

The experimental controls address a different problem: determining what
causes the observed collective differences. A field quench is compared with
matched nominal evolution so that ordinary changes over time are not mistaken
for a response to the perturbation. Persistent memory is compared with both a
memory-free condition and a scrambled-history condition, allowing us to
distinguish the effect of meaningful past interactions from the mere addition
of more prompt content. Coarse-grained block reversal measures how these
different history conditions alter the temporal ordering of the population
trajectory. Statistical inference is performed over independent
graph--environment clusters so that repeated updates within the same network
are not incorrectly treated as independent evidence. Finally, an
out-of-sample kinetic surrogate tests whether the observed collective response
can be summarized by a small number of macroscopic variables. Success would
support a useful reduced description, while failure identifies information
that those macroscopic variables omit. Together, these measurements and
controls prevent a change in a population average from being mistaken for a
complete account of collective organization, memory, and response.


The contribution is not a general-purpose agent framework, a performance
benchmark, or a model of human society.  It is a controlled finite-size study
that treats interacting LLM instances as measurable stochastic components and
examines their collective response to external fields, retained and control
histories, network organization, and coarse-graining.  The direct experiments
use two pinned instruction-model families, $N=16$ reciprocal modular networks,
and finite observation windows; the interpretation is effective and
model-specific, with no thermodynamic-limit, exact physical
entropy-production, or universality claim.  Section~\ref{sec:system-protocol}
defines the network and interventions, Section~\ref{sec:observables-inference}
defines measurement and inference, and Section~\ref{sec:evidence-hierarchy}
sets out the study roles.  Sections~\ref{sec:quench-results}--\ref{sec:memory-results}
present the quench and memory evidence, Section~\ref{sec:closure} tests reduced
closure, and Sections~\ref{sec:discussion} \& \ref{sec:conclusions} give the
interpretation and conclusions.


\section{State-separated LLM-agent networks and quench protocols}
\label{sec:system-protocol}

\subsection{Network state and update dynamics}
\label{sec:state}

For agent $i$ at attempted update $t$, the primary categorical coordinates are
belief $b_i(t)\in{-1,+1}$ and committed action
$a_i(t)\in{-1,+1}$. From a statistical-mechanical perspective, these
variables play a role similar to spin coordinates, but they represent two
distinct aspects of the agent. Belief records the agent's stated position,
whereas action records the choice that it commits to the interacting system.
Keeping them separate allows the model to represent cases in which an agent's
expressed belief and externally consequential action do not coincide.

An LLM agent also has a richer local state than a conventional spin. Its
bounded state includes confidence, commitment status, workload, a memory code,
at most three private history entries, an inbox, and an outbox. These variables
supply context for the next language-conditioned update. Consequently, two
agents with the same current belief and action can respond differently because
they have received different messages, retain different private histories, or
have different commitment states. The microscopic transition rule is therefore
not determined by the visible binary coordinates alone.

Display labels are counterbalanced across experimental conditions and decoded
back to latent spins. This prevents a particular word, label, or display
position from being permanently associated with either $-1$ or $+1$. For
statistical analysis, the selected microscopic coordinates are collected in
the projection
\begin{equation}
Y_t=(\bm b_t,\bm a_t,\bm c_t,\bm m_t).
\end{equation}
This projection provides a tractable representation of the interacting
population, but it is not the complete simulator state. In particular, private
history and message context can influence future updates even when two
projected states appear identical. Distinguishing the observed projection from
the complete state is therefore essential when interpreting memory,
Markovianity, and temporal asymmetry in the resulting population dynamics.


Let
\begin{equation}
\Xi_t=(S_{1,t},\ldots,S_{N,t},G_t,{\cal D}*t,q_t,R_t)
\end{equation}
denote all agent-private states $S*{i,t}$, graph state $G_t$, delivery operator
${\cal D}_t$, quench phase $q_t$, and the explicitly specified randomness
state $R_t$. The quench phase records whether the system is in its initial
field condition, the reversed-field condition, or the subsequent
restored-field condition. Reversing the field supplies a controlled external
perturbation, allowing us to measure how strongly the agent population
responds, while restoring it allows us to examine whether and how rapidly the
population returns toward its earlier behavior.

The scheduler is also part of the dynamical specification because it determines
which agent receives an update opportunity at each attempted update, although
it does not determine the action that the selected agent produces. Under
random-sequential updating, the order of these opportunities affects which
messages and neighbor states are available when later agents respond. Thus,
even with the same initial state and quench schedule, different scheduler
realizations can produce different microscopic trajectories. With fixed model
weights and tokenizer, reconstructed local prompts, and a defined scheduler,
$\Xi_t$ induces a time-inhomogeneous transition kernel because the external
conditions change across the quench phases. The evaluator's map
$Y_t=\phi(\Xi_t)$ deliberately omits some prompt history and randomness.
Consequently, $Y_t$ and any rolling macrostate need not be Markov.




Figure~\ref{fig:architecture} summarizes the architecture that connects local
LLM-agent interactions to the statistical-mechanical analysis. It first
distinguishes the information available to an individual agent from the
additional graph, scheduling, delivery, and randomness state maintained by the
simulator. This boundary prevents the complete simulator state from being
mistaken for information that any agent can observe or use. The figure also
shows that the scheduler only selects an agent and delivers its available
packet, whereas the LLM produces the resulting belief, action, memory update,
and outgoing signal. Finally, it follows the successive reduction from the
complete augmented state $\Xi_t$, to the evaluator's microscopic projection
$Y_t$, and then to the rolling macrostate $Z_t$. This map makes explicit where
private history and microscopic detail are omitted, providing the conceptual
basis for the later analysis of quench response, memory, and apparent
non-Markovian dynamics.

\resultfigure{figure01_architecture.pdf}{0.96\textwidth}{%
Agent information boundaries and the augmented-to-macroscopic map.  The
complete state $\Xi_t$ includes private memory, graph and schedule state; the
evaluator observes $Y_t=\phi(\Xi_t)$ and constructs rolling
$Z_t=\psi(Y_{t-w+1:t})$.  The scheduler offers a turn and delivers a packet but
does not supply the model's belief, action, memory, or signal.}{fig:architecture}

At each attempted update, a random permutation gives every agent exactly one
opportunity per sweep.  Only the scheduled instance receives its private
observation, current local state, currently delivered inbox, and, in a memory
condition, its own bounded history.  It returns a validated typed object
containing belief, action, confidence, commitment, memory state, outgoing
signal, and tool action.  One bounded repair may correct syntax; an invalid
result causes no centrally selected scientific action.  A serialized packet
is sent on one graph edge.  Fingerprints of nonrecipient peer-private states
are compared around every transition.
The corresponding information-boundary and implementation checks are described in Appendix~\ref{app:provenance}.






\subsection{History conditions and matched controls}
\label{sec:history-conditions}

The agents perform a repeated coordination task between two counterbalanced
symbolic alternatives, \texttt{amber} and \texttt{cobalt}. Neither alternative
is intrinsically preferred, and the agents are not explicitly instructed to
reach consensus. At each update opportunity, the selected agent receives a
moderate private observation favoring one alternative, its current local
state, and any packets recently delivered by its graph neighbors. It then
returns a structured decision specifying its belief, committed action,
confidence, commitment and memory states, and an outgoing signal.

Communication is therefore a form of bounded local signaling rather than
free-form conversation. A neighbor packet reports the sender's current belief,
action, outgoing signal, confidence, commitment status, and bounded-memory
category. In practical terms, it tells the receiving agent which alternative
the sender currently favors, what it acted upon, and how confidently or
firmly it holds that position. Agents do not receive their neighbors' private
observations, complete prompts, personal histories, or undelivered messages,
and they never observe the global network state.
The Markovized prompt is reconstructed from the explicitly recorded current
state and currently available packets, without carrying free conversational
history. The persistent prompt contains the same information but adds the
selected agent's last three own entries in the fixed format
\texttt{time, belief, action, memory}. Thus, two trajectories with identical
visible $(\bm b_t,\bm a_t)$ can have different conditional future laws because
their agents may remember different earlier decisions:
\begin{equation}
P(Y_{t+1}\mid Y_t,M_t)\ne P(Y_{t+1}\mid Y_t,M'_t).
\end{equation}
If $M_t$ is projected out, conditional dependence on older visible states can
remain. This provides a concrete mechanism for increased block-reversal
divergence, but it is not a claim of literal physical dissipation.

The scrambled-history control tests whether an apparent memory effect could
instead arise from the mere presence and length of an additional prompt
section. It retains the same agent's past opportunity times and uses the same
number of entries, field names, and format, while the belief, action, and
memory content is deterministically randomized from a prespecified control
seed. No donor agent, future time, or hidden environmental outcome is used.
Token counts and their paired differences are measured, although exact
turn by turn equality is not assumed. Prompt balance diagnostics are shown in
Figure~\ref{fig:prompt-balance} of Appendix~\ref{app:sensitivities}.




\subsection{Field quench and restoration protocol}
\label{sec:quench-protocol}

Each trajectory contains 15 baseline sweeps, 15 intervention sweeps, and 15
restoration sweeps. One sweep consists of $N=16$ attempted updates, giving each
agent one update opportunity on average under random-sequential scheduling.
The baseline allows collective behavior to develop under the initial fields;
the intervention permits the reversed evidence to propagate through repeated
local interactions; and the equally long restoration phase provides a
comparable window for observing recovery. These phase lengths were fixed
prospectively rather than selected after inspecting the trajectories. They
define a finite observation horizon and do not assume that stationarity is
reached.

In each field-quench arm, the balanced vector of private fields changes sign at
the first boundary. An agent previously receiving evidence favoring
\texttt{amber}, for example, then receives evidence favoring
\texttt{cobalt}, while the graph, model, and interaction rules remain
unchanged. The agents are not told that a global intervention has occurred, so
the collective response must develop from altered local evidence and
neighbor-to-neighbor signaling. Matched nominal trajectories retain the
initial field and control for changes that would occur without the quench. At
the second boundary, the original fields are restored, allowing the experiment
to measure whether collective organization returns, remains delayed, or
persists over the available horizon. This is the logic of a finite-system
quench and counter-quench, without assuming quantum dynamics or a
thermodynamic limit.

Each model family contributes six independent graph/environment clusters. A
cluster is a matched experimental replication block, not a community within
the modular graph. It combines one graph realization with its private-field
pattern, initial states, label mapping, and associated randomization. Each
cluster contains four matched trajectories: nominal Markovized, field-quench
Markovized, field-quench persistent memory, and field-quench scrambled
history. All use $N=16$, a reciprocal modular fixed-degree graph, coupling
$J=0.8$, and balanced disordered initialization. The modular topology provides
local and community pathways through which signals can propagate, while
matching graphs, fields, initial states, label maps, update orders, and other
random variates reduces irrelevant between-condition variation. Because the
four trajectories within a cluster share these design elements, the
graph/environment cluster, rather than an individual agent or update, is the
unit of statistical evidence.




\section{Statistical-mechanical observables and\texorpdfstring{\\}{ }inferential design}
\label{sec:observables-inference}

\subsection{Macroscopic and information-theoretic observables}
\label{sec:order-compatibility}

Belief and action order and their overlap are,
\begin{equation}
m_b=\frac{1}{N}\sum_i b_i,\qquad
m_a=\frac{1}{N}\sum_i a_i,\qquad
q_{ba}=\frac{1}{N}\sum_i b_i a_i.
\end{equation}
Disagreement is the fraction of connected pairs with unequal states.  For a
symmetric comparison layer $W_s=(W+W^\top)/2$, we define
\begin{equation}
H_{\mathrm{ref}}=-\frac{J_b}{2}\bm b^\top W_s\bm b
-\frac{J_a}{2}\bm a^\top D_s\bm a
-K\bm a^\top\bm b-\bm h^\top\bm b-\bm g^\top\bm a .
\label{eq:href}
\end{equation}
We report $e_{\mathrm{ref}}=H_{\mathrm{ref}}/N$.  Equation~\eqref{eq:href} is
an effective compatibility coordinate anchored in a
symmetric reference model.  We do not infer a Gibbs law for the LLM process
and do not call this literal energy.  Likewise
$N\operatorname{Var}(e_{\mathrm{ref}})$ is an energy-fluctuation observable,
not heat capacity unless an independently justified effective temperature
exists.
The belief susceptibility is reported without physical-temperature
normalization,
\begin{equation}
\chi_b=N\operatorname{Var}(m_b),
\end{equation}
within a prescribed rolling window.

For discretized collective words $x$ in a rolling window, Shannon
configuration
entropy is
\begin{equation}
S_{\mathrm{config}}=-\sum_x \hat p(x)\log\hat p(x).
\end{equation}
This is the usually used Shannon functional \cite{shannon1948}.
The entropy rate is a bounded-history plugged in estimate of the next word
unpredictability.  Single-variable marginal entropy measures local occupancy,
pairwise mutual information measures dependence between variables.  Total
correlation \cite{watanabe1960},
\begin{equation}
{\cal T}=\sum_{k=1}^{2N}H(X_k)-H(X_1,\ldots,X_{2N}),
\end{equation}
separates marginal uncertainty from synchronous joint dependence.

The plug in joint entropy of 32 binary variables in a short window has a large
finite-sample floor \cite{paninski2003}.  We therefore retain the raw estimate
but independently circular-shift each variable 200 times.  Shifts preserve
each marginal occupancy and within-variable temporal ordering while breaking the
synchronous cross variable relations.  We report raw, null mean, their
untruncated difference, and a normalized difference divided by the sum of
marginal entropies.  The same audit is applied to pairwise and edge mutual
information.  A negative adjusted estimate remains negative.
The primary five-sweep representation is
\begin{align}
Z_t=(&m_b,m_a,q_{ba},e_{\mathrm{ref}},v_E,S_{\mathrm{config}},h_{\mathrm{rate}},
{\cal T}_{\mathrm{adj}},I_{\mathrm{pair,adj}},I_{\mathrm{edge,adj}},\nonumber\\
&\chi_b,\rho_b,D_b)_t .
\end{align}
The three and seven sweep windows are prespecified sensitivities.  This is an
observable reduced representation, not a complete thermodynamic state.







\subsection{Temporal asymmetry, spatial correlation, and persistence diagnostics}
\label{sec:path-asymmetry}

For a discrete observed word sequence and block length $L$, let
$\hat P_L(x_0,\ldots,x_L)$ be the regularized empirical forward block law.
We compute
\begin{equation}
{\cal I}_L=\frac{1}{L}
\KL\left[\hat P_L(x_0,\ldots,x_L)\parallel
\hat P_L(x_L,\ldots,x_0)\right]
\label{eq:pathkl}
\end{equation}
in nats per attempted update.  The primary $L=3$ estimator uses pseudocount
0.5.  Values at $L\in\{2,4\}$ and pseudocounts 0.1 and 1.0 are sensitivities.
Five hundred time permutations provide a panel-specific finite-length floor,
and the adjusted statistic is the observed value in
equation~\eqref{eq:pathkl} minus its mean floor.
Matched arms share the same shuffle tape.  The machine-readable sensitivity
table retains both the empirical permutation-null interval and the Monte Carlo
standard error of its mean.  These quantify different objects, dispersion of
the null statistic and numerical uncertainty in the estimated bias floor,
respectively.  The added uncertainty columns reproduce, rather than replace,
the prespecified floor mean and adjusted contrast.

For a fully observed stationary Markov chain with transition kernel $P$ and
stationary law $\pi$, the current
\begin{equation}
J_{xy}=\pi_xP_{xy}-\pi_yP_{yx}
\end{equation}
can define an exact Markov entropy-production rate under the usual support
conditions.  We do not apply that interpretation to projected LLM trajectories.
Persistent histories, semantic content, and bounded observations can violate
first-order closure.  Equation~\eqref{eq:pathkl} is instead a coarse-grained
time-reversal diagnostic and, under additional assumptions, may be a lower
bound on hidden irreversibility \cite{roldan2012,busiello2019,liang2023}.

Three secondary observables resolve
spatial organization, temporal persistence, and the shape of the finite-size
order-parameter distribution.  First, for a phase window $W$ we define the
connected belief correlation at graph distance $d$,
\begin{equation}
C_b(d;W)=\frac{1}{|{\cal P}_d|}\sum_{(i,j)\in{\cal P}_d}
\left[\langle b_i b_j\rangle_W-\langle b_i\rangle_W
\,\langle b_j\rangle_W\right],\qquad
{\cal P}_d=\{(i,j):i<j,\ d_G(i,j)=d\}.
\label{eq:connected-correlation}
\end{equation}
Shortest paths use the actual reciprocal modular delivery support.  The
bracket averages post-update observations at attempted-update resolution
within one phase of one
trajectory, node pairs are averaged inside that trajectory and never counted
as independent replicates.
Second, the normalized magnetization autocorrelation and its fixed truncation
are
\begin{align}
\rho_m(\ell)&=\frac{\langle[m_b(t)-\bar m_b]
\,[m_b(t+\ell)-\bar m_b]\rangle}
{\langle[m_b(t)-\bar m_b]^2\rangle},\nonumber\\
\tau_{\mathrm{int}}^{(\ell_{\max})}&=\frac{1}{2}+
\sum_{\ell=1}^{\ell_{\max}}\rho_m(\ell).
\label{eq:autocorrelation}
\end{align}
The primary descriptive truncation is $\ell_{\max}=2N=32$ attempted updates,
$N$ and $3N$ are fixed sensitivities.  Zero-variance phase series are left
undefined rather than assigned artificial persistence.  The quantity is a
finite-window correlation sum reported in attempted-update units, not evidence
of critical slowing down or an effective sample-size estimate.  Its absolute
shorter lag scale includes persistence induced by changing at most one agent on
each random-sequential update; matched arms share that update convention.
Quench and restoration windows can also be nonstationary, so their sums are
finite-window persistence diagnostics rather than equilibrium integrated
correlation times.

Third, we calculate the Binder cumulant
\begin{equation}
U_4=1-\frac{\langle m_b^4\rangle}
{3\langle m_b^2\rangle^2}
\label{eq:binder}
\end{equation}
separately for each phase and trajectory \cite{binder1981}.  Near-zero second moments are
reported as undefined.  With only $N=16$, equation~\eqref{eq:binder} is a
finite-size distribution-shape diagnostic; no crossing, critical point, or
thermodynamic-limit transition is inferred.  A post-reconstruction estimator
sensitivity repeats the calculation over each full phase and its early and
late halves.  It compares the mean of the six trajectory-level cumulants with
a pooled-moment construction; both uncertainty calculations resample complete
graph--environment clusters, never updates.





\subsection{Recovery estimands and exact inference}
\label{sec:inferential-design}

For each hold out graph environment cluster and model, the nominal manifold is
fitted only from nominal trajectories in the other five clusters.  Training
medians impute nonfinite coordinates; training means and scales standardize
features, shrinkage covariance with a prespecified ridge fraction defines a
Mahalanobis like distance.  The nominal 95th percentile threshold is computed
from those training distances.  The hold out cluster contributes neither a
scale nor a threshold.

We summarize peak departure, integrated response, path length, final residual,
and threshold re-entry.  The prospective cross model recovery contrast is
\begin{equation}
\Delta d_{\mathrm{rec}}=
\overline d_{31:35}-\overline d_{41:45},
\end{equation}
which can be positive, zero, or negative.  It replaces a historical maximum
minus final window quantity whose sign was structurally constrained.
A complete graph/environment trajectory cluster is the
inferential unit.

The hypotheses were specified before the cross-model outcomes were examined.
H1 compares Granite field quench peak
departure with matched nominal evolution at a one-sided $\alpha=0.02$.  H2
compares persistent with Markovized adjusted reversal divergence.  H3 compares
persistent with scrambled history.  H4 tests the fixed recovery contrast.
H2 to H4 form a one sided Holm family with total $\alpha=0.03$.  Exact cluster
sign flips are the primary tests, deterministic 10,000 cluster bootstraps give
intervals.  The sign flip enumeration is exact conditional on the observed
paired magnitudes under a sign-symmetry null; it is not described as treatment label randomization.
Agent and time-point counts do not increase the inferential sample size.
The exact test, multiplicity, and cluster bootstrap details are retained in Appendix~\ref{app:study-inference}.




\section{Study design and evidential roles}
\label{sec:evidence-hierarchy}

\subsection{Exploratory discovery and corrected earlier evidence}
\label{sec:memory-study-roles}

The memory evidence developed in three stages, an exploratory experiment, a
subsequent prospective replication, and the present cross-model extension with
an additional history control.

The evidence for a memory effect was accumulated in distinct stages with
different inferential roles. The exploratory stage identified increased
temporal asymmetry under persistent history as a candidate phenomenon, while
the prospective replication tested a contrast specified after that discovery.
The later cross-model experiment then examined the effect under the frozen
design and matched Markovized and scrambled-history controls for Granite and
Qwen.
Figure~\ref{fig:memory} presents these stages separately, with each point
representing the effect observed within one independent graph/environment
cluster and intervals constructed by resampling complete clusters. This
separation allows the reader to assess consistency across stages and
heterogeneity between model families without allowing the hypothesis-generating
or earlier replication data to narrow the cross-model confirmatory interval
retroactively.

\resultfigure{figure02_memory_evidence_stages.pdf}{0.94\textwidth}{%
Exploratory discovery, prospective replication, and cross-model memory contrasts are
shown separately.  Points are graph/environment-cluster effects; intervals
resample complete clusters.  The exploratory and replication-stage estimates
are not retroactively pooled into the cross-model confirmatory interval.}{fig:memory}

The audit of the earlier quench experiment preserves every recorded prompt,
completion, categorical choice, and trajectory.  The original analysis
computed a field panel's recovery
threshold from its own baseline, despite protocol text specifying a
training-only threshold.  The corrected analysis passes an explicit
cluster-to-training-threshold map from the 'leave one cluster out' fitting into the
recovery summaries.

The earlier study's third hypothesis used the maximum recovery distance minus
the final-five mean.  That value remains numerically reproducible, but the
construction is
nonnegative for nearly every nonconstant trajectory.  Its raw and Holm-adjusted
sign-flip probabilities are retained as historical audit values and removed
from inferential support.  The corrected report distinguishes
whether the numerical criterion was met, whether a directional test is valid,
and whether the time-resolved path is consistent with return.

Figure~\ref{fig:v14series} displays the corrected time series from the earlier
quench experiment and places the nominal and field-reversal trajectories on a
common time axis. This comparison distinguishes changes associated with the
reversed private fields from background evolution that also occurs without an
intervention. The yellow interval identifies the period during which the
fields are reversed, while the dotted boundary marks their restoration and
therefore separates the response phase from the recovery phase. The displayed
observables provide complementary views of the dynamics: effective reference
energy summarizes interaction compatibility, entropy measures configuration
diversity, and distance quantifies displacement from nominal behavior in the
cluster-excluded geometry. Together, the panels show whether the perturbation
produces a coordinated departure during the quench and whether that departure
contracts, persists, or changes form after restoration. The correction applies
to the derived analysis used to construct these summaries; the underlying
agent trajectories are unchanged.
Cluster-level recovery and delayed-audit material is retained in Appendix~\ref{app:study-inference}, including Figures~\ref{fig:v14recovery} and~\ref{fig:v14audit}.

\resultfigure{figure03_corrected_quench_time_series.pdf}{0.94\textwidth}{%
Corrected time series from the earlier experiment for nominal and field-reversal trajectories.  The
yellow interval marks the quench; the dotted boundary marks restoration.
Reference energy is effective, entropy is configuration diversity, and
distance is measured in the cluster-excluded nominal geometry.}{fig:v14series}




\subsection{Cross-model confirmation and control logic}
\label{sec:cross-model-design}

To test whether the observed behavior extended beyond the model family in
which it was first identified, the cross-model study used
Qwen/Qwen2.5-7B-Instruct as the original family and
IBM Granite-3.3-8B-Instruct as an independently developed comparison family
\cite{qwen2024,granite33model}. Both models were pinned to immutable revisions,
with the full revision identifiers recorded in Appendix D and the repository
manifests. Using two separately developed instruction-tuned families provides a
stronger test of generality than repeating the experiment with another
checkpoint from the same family. At the same time, the inference settings were
held fixed: both models used the same pinned software stack, NF4 double
quantization, BF16 computation, decoding temperature 0.5, top-$p=0.9$, and a
maximum of 96 generated tokens. The decoding temperature is a stochastic
text-generation setting and is not interpreted as a physical or inferred
thermodynamic temperature.

Before outcome evaluation, a limited engineering pilot verified that both
families could execute the same protocol reliably. These checks covered model
loading, JSON-schema validity, state occupancy, both transition directions,
information isolation, packet delivery, prompt-token balance, latency, and
projected compute requirements. The pilot did not calculate quench, memory, or
block-reversal contrasts and therefore supplied no outcome evidence. The
prospectively specified experiment then comprised 48 trajectories and 34,560
attempted local decisions across the two model families. Every complete
trajectory was retained, including trajectories with weak, null, or
unfavorable effects, so inclusion in the confirmatory analysis did not depend
on the observed result.






\section{Quench response, restoration, and finite-size organization}
\label{sec:quench-results}

\subsection{Cross-model response and recovery}
\label{sec:field-response}

The prespecified quench-response test (H1) was supported in the independent
Granite family. The mean field-minus-nominal peak-distance contrast was
42.263 within-model distance units, with a 95\% interval of
$[24.316,59.303]$. A positive contrast means that reversing the private fields
drove the population farther from its nominal trajectory than would be
expected from unperturbed variation. All six Granite cluster effects were
positive. The separately retained Qwen results showed the same direction in
all six clusters, with a mean field-minus-nominal peak contrast of 112.691
units. Complete cluster results were retained regardless of their direction or
magnitude.

Figure~\ref{fig:crossquench} shows the nominal and field quench trajectories
separately for the two model families. Within each panel, the figure shows when
the quenched population departs from its matched nominal evolution and how that
separation changes after the original fields are restored. The vertical scales
should not be compared directly between Qwen and Granite because the nominal
geometry, covariance, scaling, and thresholds are estimated separately within
each model. These quantities are fitted while excluding the displayed cluster,
so a cluster does not define the geometry in which its own response is
measured.

\resultfigure{figure04_cross_model_quench.pdf}{0.94\textwidth}{%
Model-specific field-quench and nominal trajectories.  Curves are means over
six graph/environment clusters and bands show between-cluster dispersion.
Nominal fits, scaling, covariance, and thresholds exclude the displayed
cluster.}{fig:crossquench}

The prespecified restoration test (H4) was also supported. The fixed
 restoration contrast was 52.541 distance units, with a 95\%
interval of $[35.406,68.673]$. Its positive value indicates that distance from
the nominal regime declined between the early and late restoration windows,
and this decline occurred in every model to cluster pair. Unlike the historical
recovery statistic used in the earlier experiment, this contrast was defined
in advance and permits either sign.

Movement toward the nominal regime does not necessarily imply complete return
within the finite observation horizon. Every Qwen trajectory re-entered its
threshold, estimated from nominal training data, six sweeps after restoration.
Only one of the six Granite trajectories re-entered by sweep 45; the remaining
five had final five sweep means above their corresponding thresholds. Granite
therefore showed clear movement toward restoration together with slower or
incomplete finite-horizon recovery. The quench and counter quench paths need
not retrace one another because messages, actions, and histories accumulated
during the intervention alter the state from which restoration begins. These
recovery measurements are finite time trajectory descriptors, not claims of
equilibrium relaxation.






\subsection{Spatial and finite-window organization}
\label{sec:spatial-results}

The preceding response measures establish whether the population moves away
from nominal behavior, but they do not show how that change is organized
across the network or how long the resulting collective state persists.
Connected graph-distance correlations address the first question by measuring
whether neighboring agents reorganize together and whether that organization
extends across the predefined communities. Finite-window autocorrelation and
Binder statistics address the second by describing persistence and the shape
of the observed order-parameter distribution. These are secondary descriptive
analyses of the same recorded trajectories: they introduce no additional
simulation arms and do not modify the confirmatory estimands. Because they
were specified during reconstruction after the confirmatory aggregate results
were known, they are not treated as a new prospective hypothesis family.

Connected correlation matrices from
equation~\eqref{eq:connected-correlation} are calculated separately within
every complete model--graph--environment trajectory and phase. Their
graph-distance profiles are then summarized over the six independent
clusters. Figure~\ref{fig:correlations} shows the cluster-averaged matrices for
the persistent-history quench window together with phase-resolved correlation
profiles by graph distance. The community boundaries are fixed by the graph
construction rather than selected from the observed correlations, so the
analysis asks how behavior aligns with a predefined interaction structure.

At graph distance one, the persistent-history quench-minus-baseline change in
connected correlation was $0.00238$, with a 95\% interval of
$[-0.01750,0.02277]$, for Qwen and $-0.02032$, with a 95\% interval of
$[-0.05486,0.00947]$, for Granite. Both intervals include zero, indicating
that the available clusters do not resolve a consistent nearest-neighbor
change during the quench. These intervals describe finite-window spatial
reorganization and are not assigned a confirmatory disposition.

\resultfigure{figure05_graph_distance_correlations.pdf}{0.94\textwidth}{%
Connected belief correlations in the persistent-history arm.  The upper
panels average complete $16\times16$ connected-correlation matrices over all
six clusters in each model; black rules mark the two predefined communities.
Agent indices align across graph realizations only through those community
memberships, so fine node-level texture in an averaged matrix has no
cross-realization identity interpretation.
The lower panels show equation~\eqref{eq:connected-correlation} by graph
distance, with 95\% complete-cluster bootstrap intervals.  Individual pairs
and updates do not enter the inferential sample size.}{fig:correlations}

The finite-window diagnostics examine temporal organization from a
complementary perspective. Appendix~\ref{app:persistence-binder} and
figure~\ref{fig:persistence-binder} combine the full autocorrelation curve with
the fixed-lag integral in equation~\eqref{eq:autocorrelation}, the Binder
statistic in equation~\eqref{eq:binder}, and empirical magnetization
occupancy. Showing occupancy alongside the Binder statistic is important
because a finite-window Binder value can be strongly affected by the limited
distribution of observed magnetization states. Autocorrelation also measures
a different property from block-reversal divergence: autocorrelation is a
time-symmetric measure of two-point persistence, whereas block reversal
compares the frequencies of forward and reversed path words.

At the fixed two-sweep truncation, the persistent-minus-Markovized recovery
change in $\tau_{\mathrm{int}}$ was $-2.46995$ updates, with a 95\% interval of
$[-5.79214,0.17724]$, for Qwen and $-4.14109$ updates, with a 95\% interval of
$[-19.88286,11.60067]$, for Granite. The field-Markovized
quench-minus-baseline Binder change was $-4.80504$, with a 95\% interval of
$[-7.42824,-2.20832]$, for Qwen and $0.16643$, with a 95\% interval of
$[-0.09447,0.52063]$, for Granite. The accompanying source table retains
full-window, early-half, and late-half Binder estimates under both
cluster-first and pooled-moment constructions. Differences among these
estimates are used to diagnose finite-window and aggregation sensitivity, not
to select the construction that produces the largest separation. All results
are reported regardless of sign or compatibility with zero, and none is
interpreted as evidence of an equilibrium correlation time, critical slowing
down, a Binder crossing, or a phase transition.







\section{Memory and projected temporal asymmetry}
\label{sec:memory-results}


\subsection{Persistent, Markovized, and scrambled histories}
\label{sec:persistent-markovized}

The three history conditions distinguish between having no explicit personal
history, having genuine history, and merely having additional history-shaped
text. In the Markovized condition, an agent sees its current state but not its
earlier decisions. In the persistent condition, it also sees its last three
genuine belief, action, and memory records. The scrambled condition provides a
similarly formatted history section, but its contents are randomized and no
longer correspond to the agent's actual past.

Block-reversal divergence measures whether short sequences of projected
population states occur differently from the same sequences read in reverse.
A larger value indicates a stronger temporal direction in this reduced
description of the trajectory. Relative to the Markovized condition, genuine
history increased the bias-adjusted block-reversal divergence by
$0.05438$ nats per attempted update, with a 95\% interval of
$[0.03459,0.07548]$. This prespecified comparison, H2, was supported. It shows
that providing an agent with its actual recent history changes the temporal
organization of the observed population dynamics beyond what is produced by a
current-state prompt alone.

The comparison with scrambled history provides a stricter control. Genuine
history increased adjusted block-reversal divergence by $0.04845$ nats per
attempted update relative to the scrambled condition, with a 95\% interval of
$[0.02190,0.07526]$. This prespecified comparison, H3, was also supported.
Because both prompts contain similarly formatted history sections, the result
indicates that the relationship between the displayed history and the agent's
actual past matters beyond the mere presence or approximate length of
additional prompt text.

Some adjusted values for individual arms are negative because the correction
subtracts the average divergence expected from the finite-sample shuffle
baseline. An observed value below that baseline therefore produces a negative
adjusted score. This does not represent negative physical entropy production,
and the values are not truncated to zero before computing the paired H2 and H3
contrasts. The positive contrasts show a difference between history
conditions, not that either arm has positive absolute thermodynamic entropy
production.

Figure~\ref{fig:v15memory} displays the paired cluster-level contrasts for the
two model families. Each point represents one independent
graph/environment-cluster comparison, so its position relative to zero shows
whether genuine history increased or decreased adjusted block-reversal
divergence within that matched cluster. The separate model summaries also make
visible whether the size and consistency of the effect differ between Qwen and
Granite.

\resultfigure{figure06_memory_controls.pdf}{0.94\textwidth}{%
Per-cluster genuine-memory contrasts for both pinned model families.  Short
horizontal segments indicate model means.  The scrambled-history arm controls
prompt section and approximate length while destroying temporal relation to
the current trajectory.}{fig:v15memory}

No exact thermodynamic interpretation follows from these results. The block
words summarize selected categorical coordinates and omit the natural-language
messages and part of each agent's private state. Sensitivity tables for block
length, pseudocount, conditional memory depth, state occupancy, and prompt
balance test how strongly the result depends on these analytical choices. The
prompt-balance diagnostic in figure~\ref{fig:prompt-balance} of
Appendix~\ref{app:sensitivities} records the scope of the format and
approximate-length control.






\subsection{Model heterogeneity and estimator dependence}
\label{sec:memory-heterogeneity}

The pooled H2 and H3 results show that genuine history changes the projected
dynamics on average across the prespecified model stratified clusters. A
supported pooled effect, however, does not imply that Qwen and Granite respond
to history in the same way or with the same magnitude. It also does not imply
that every reasonable construction of the block reversal estimator produces
an identical result. We therefore examine model heterogeneity and estimator
sensitivity separately to define the scope of the supported memory findings.

For persistent relative to Markovized prompting, the model-specific mean
contrast was $0.03041$ nats per attempted update for Qwen, with five of six
cluster effects positive, and $0.07835$ for Granite, with all six positive.
Both families therefore showed the same average direction for H2, although the
effect was larger and more consistent in Granite. For persistent relative to
scrambled history, the mean contrast was $0.00689$ for Qwen, with four of six
cluster effects positive, and $0.09001$ for Granite, with all six positive.
The separate Qwen scrambled-control sign-flip value was $0.21875$ and was not
itself a confirmatory test. The prespecified pooled H3 remained supported
because its inferential units were the 'model stratified' cluster pairs. The
appropriate conclusion is therefore that temporally related history has a
supported pooled effect, with a strong Granite contribution, rather than that
the content-specific effect is uniform across model families.

The magnitude of the memory contrast also depends on how block-reversal
divergence is estimated. Block length determines how much temporal structure
is represented, the pseudocount controls sparse word frequencies, and the
shuffle floor corrects finite-sample bias. The pooled memory contrasts remain
positive across most prespecified combinations of these choices, but their
magnitudes vary substantially. At block length four and pseudocount one, for
example, the persistent-minus-scrambled contrast falls to $0.00024$ nats per
attempted update, and the Granite component becomes slightly negative.
Conditional-memory information is also not consistently larger in the
persistent condition under every construction.

These sensitivity results do not reverse the prespecified H2 or H3
dispositions. They limit the claim to the prespecified coarse graining,
bias-floor correction, and paired controls, rather than supporting an
estimator-independent entropy-production rate. The block-length, pseudocount,
and shuffled-floor results are displayed in
Figure~\ref{fig:path-sensitivity} of Appendix~\ref{app:sensitivities}.

Taken together, H1 to H4 addresses four distinct questions. H1 tests whether the
field reversal drives the system beyond nominal evolution, H2 tests whether
genuine history matters relative to current-state prompting, H3 tests whether
temporally related history matters beyond a format-matched scrambled control;
and H4 tests movement toward the nominal regime during the fixed restoration
window. Their joint support is therefore a synthesis of complementary results,
not evidence for a single universal effect.


Figure~\ref{fig:effects} brings the four prespecified tests together as a
summary of the main confirmatory findings. These tests address distinct
questions rather than repeated measurements of one underlying effect. H1 asks
whether the field quench produces a departure beyond nominal evolution, while
H4 measures movement back toward the nominal regime between the early and late
restoration windows. H2 compares genuine persistent history with
current-state-only Markovized prompting, and H3 compares genuine history with
the format-matched scrambled-history control.
The left panel displays H1 and H4 in within-model distance units, whereas the
right panel displays H2 and H3 in nats per attempted update. In each case, the
point is the estimated effect, the horizontal bar is its 95\%
complete-cluster bootstrap interval, and the dashed vertical line marks zero.
All four intervals lie on the positive side of zero, but positivity has a
different meaning for each test: greater quench-induced departure for H1,
greater early-to-late restoration for H4, and greater adjusted block-reversal
divergence under genuine history for H2 and H3. The separate panels prevent
distance and information quantities from being compared as though they shared
a common scale or effect magnitude.
Under the prespecified tests, exact cluster sign-flip calculations, and
multiplicity allocation, all four criteria were met. The figure therefore
provides a compact synthesis of the supported outcomes, but it does not combine
them into a single composite effect. The model-specific differences,
finite-window interpretation, and estimator limitations described above remain
part of the result.

\resultfigure{figure07_confirmatory_effects.pdf}{0.94\textwidth}{%
Four prespecified effects with 95\% complete-cluster bootstrap intervals.
The left panel shows the field-quench response (H1) and fixed-window
restoration contrast (H4) in within-model distance units.  The right panel
shows genuine persistent history relative to Markovized prompting (H2) and
scrambled history (H3) in nats per attempted update.  Points are estimated
effects, horizontal bars are bootstrap intervals, and dashed vertical lines
mark zero.  Distance and information quantities are kept on separate axes
because their numerical magnitudes are not comparable.  Prespecified
dispositions use exact cluster sign flips and the multiplicity allocation.}{fig:effects}





\section{Reduced dynamical descriptions and closure limits}
\label{sec:closure}


\subsection{Information retained beyond mean order and the kinetic surrogate}
\label{sec:representation}

Magnetization records directional order but can remain near zero when half the
population reorganizes in opposing communities.  Configuration entropy then
records collective state diversity; marginal entropy records local occupancy;
total correlation distinguishes independent uncertainty from coordinated
dependence; effective energy records compatibility with the symmetric
reference; susceptibility and energy variance record fluctuations;
autocorrelation records finite-window persistence; and path reversal records
temporal ordering.  These
quantities are complementary only when their estimators and scales are kept
distinct.

The representation analysis of the earlier quench experiment uses transparent multinomial logistic
regression with leave-one-cluster-out preprocessing.  The delayed permutation
audit tests both absolute balanced accuracy and the increment over order-only
features.  This is not a claim that entropy outperforms all ordinary summaries.
It tests whether a predefined collective representation retains condition
information discarded by a reduced representation.  The field quench
trajectory supplies a mechanistic example, energy, dependence, entropy rate,
and fluctuation can change while mean magnetization remains balanced.

We fit logistic belief and action responses only to the immutable
microscopic response table from the prospective replication stage.  The belief model includes private field, delivered
neighbor field times controlled coupling, previous belief, and previous action;
the action model includes updated belief and previous action.  No coefficient
is refitted to either later quench trajectory set.  The surrogate is run on the
same graph generator, initial state generator, asynchronous update tape, field
schedule, and intervention times. It is a deliberately low-dimensional
kinetic Ising comparison rather than a fitted description of every asymmetric
or history-dependent response \cite{aguilera2021}.

Figure~\ref{fig:surrogate} compares the direct Qwen-agent trajectories from the earlier field-quench experiment with trajectories generated by the low-dimensional kinetic surrogate. The solid blue curves represent the observed LLM-agent dynamics, whereas the dashed orange curves show what the surrogate predicts when it is run with the same graph, initial conditions, update schedule, and field changes. This comparison is genuinely out-of-sample: none of the surrogate coefficients are fitted or adjusted using the quench trajectories displayed in the figure. The eight panels examine different aspects of the collective response, including belief and action magnetization, belief-action overlap, effective reference energy, configuration entropy, susceptibility, finite-window correlation time, and shared-response distance. The yellow region marks the field-reversal interval, with its boundaries identifying the quench and restoration events. A successful reduced description would not need to reproduce every microscopic fluctuation, but it should capture the direction, timing, relative amplitude, and recovery of the main collective changes across these observables. Agreement therefore identifies behavior that can be explained by the surrogate's small set of local variables, while disagreement reveals collective, semantic, asymmetric, or history-dependent information that the reduced closure omits.

\resultfigure{figure08_direct_surrogate_quench.pdf}{0.94\textwidth}{%
The out of sample field-quench comparison between direct Qwen trajectories and
the replication stage kinetic surrogate.  Matching direction, peak timing, normalized
amplitude, recovery, and route area are assessed separately.  Surrogate
failure is evidence about the limits of the closure, not a failed LLM
experiment.}{fig:surrogate}

In the five coordinate comparison geometry, the direct field minus nominal
peak response difference is only 0.142 units on average and is nonzero in one
of six clusters, whereas the surrogate difference is 1.467 units and is
positive in four clusters.  The corresponding field-minus-nominal
energy entropy route-area differences are 0.043 and 0.695.  The direct shared
response peaks at the quench boundary in one cluster and is otherwise dominated
by initial relaxation; the surrogate places its peak at the quench or
counter-quench boundary in four clusters.  Both descriptions return close to
their own final shared-response baseline, but the surrogate neither reproduces
the direct amplitude nor its timing.  This also explains why the richer
macrostate distance used in the earlier quench analysis can show a large quench pulse that a short kinetic closure
over a small shared coordinate set misses.




\subsection{Captured response, failure modes, and finite-size context}
\label{sec:closure-limits}

This comparison identifies three possible failure modes.  First, a homogeneous
logistic neighbor coefficient cannot capture state-dependent semantic use.
Second, action and belief persistence can introduce lags not represented by a
single synchronous mean field.  Third, hidden memory and workload can produce
different future laws at the same projected spin configuration.  Adding
unconstrained parameters after observing quench trajectories would obscure
these limits, so the closure remains low-dimensional and out of sample.

CPU trajectories at $N\in\{8,16,32,64\}$ supply a dense effective-model size
context.  They are not direct LLM finite-size scaling and cannot establish a
critical exponent or thermodynamic-limit phase transition.  The direct
$N=16$ paths remain empirical anchors.

The effective-model size context is displayed in figure~\ref{fig:surrogate-size}
of Appendix~\ref{app:sensitivities}.




\section{Discussion}
\label{sec:discussion}

Treating the LLM-agent network as a stochastic interacting system begins with
implemented rather than metaphorical objects, agent instances receive local
prompts, exchange packets, update on a recorded schedule, and change
categorical state.  Magnetization, Shannon entropy, mutual information, total
correlation, and path divergence are then mathematical observables of those
records.  State separation is central to this interpretation because it
distinguishes the complete augmented state that determines a transition from
the local information available to an agent and from the lower-dimensional
projection available for population-level measurement.

The quench experiments use this formulation to measure finite-horizon response
and recovery rather than equilibrium relaxation.  Both model families depart
from their within-model nominal geometries after field reversal, and the
fixed-window restoration contrast supports motion back toward the restored
nominal regime.  The difference between complete Qwen threshold re-entry and
slower or incomplete Granite re-entry within the observation horizon remains
important: the common directional response does not imply a common recovery
time or complete return.  Recovery times and threshold crossings are
trajectory descriptors, and the quench and counter-quench paths need not
retrace because messages, actions, and history change the state reached at
restoration.

The history interventions clarify what memory means under coarse-grained
observation.  Persistent history differs from an explicit current-state
prompt, and comparison with scrambled history asks whether temporally related
content matters beyond prompt format and approximate length.  The distinction
between complete augmented state, observable projection, and rolling
macrostate explains how private memory can generate history dependence in the
projected process.  Positive paired contrasts nevertheless do not establish
positive absolute entropy production: some arm level bias-adjusted values are
negative, the observed path words omit natural language content and private
state, and the measured reversal divergence is a coarse-grained projected
statistic rather than an exact physical entropy production rate.

Model heterogeneity and negative findings limit generalization.  The
persistent versus Markovized contrast has the same direction across the two
families, whereas the content specific scrambled history contrast is much
weaker and less uniform for Qwen than for Granite.  Estimator choices also
change the magnitude of the memory contrasts.  In the exploratory evidence,
degree  and traffic matched nonreciprocity did not produce a reliable increase
in irreversibility, and partition and message corruption trajectories in the
earlier quench experiment remained close to nominal under their tested
construction.  These results prevent the broader assertion that any directed
graph, entropy statistic, or history prompt produces the same collective
effect.

The secondary spatial, persistence, Binder, and closure analyses add distinct
but limited information.  Connected correlations describe how finite-window
organization changes across the fixed modular graph, while autocorrelation
summarizes persistence of the projected collective coordinate.  Binder values
are shown with occupancy and aggregation sensitivities,  they do not establish
a crossing, critical slowing down, or a phase transition.  The out of sample
kinetic surrogate captures some direction and recovery features but misses
parts of the direct response amplitude and timing.  Its inexpensive size
context is not direct LLM finite size scaling, so surrogate failure identifies
limits of the reduced closure rather than a failed agent experiment.

These measurements support an effective statistical mechanical description,
not a literal identification of the network with a thermodynamic material.
The reference energy is effective because the LLM transition rule is not
derived from its Hamiltonian, and decoding temperature is a generation setting
rather than an independently inferred thermodynamic temperature.  No
free energy like coordinate is emphasized because such a temperature is
unavailable.  The value of the framework is instead that it keeps collective
organization, fluctuations, dependence, response, and temporal ordering
operationally distinct while exposing the assumptions behind each reduction.

Several limitations bound the scope.  Two 7B instruction model families do
not establish that it will apply identically to every model,
and the direct systems contain only $N=16$ agents,
so no thermodynamic limit transition is claimed.  The binary projection
compresses semantic messages and private state.  The confirmatory study uses
one reciprocal modular topology and one coupling/noise anchor.  The
scrambled history placebo matches format and approximate token distribution
but cannot make semantic content exchangeable turn by turn.  Short window
entropy, dependence, autocorrelation, and path estimates remain sensitive to
coarse graining and finite sampling even after their null floor and
aggregation checks.

Further work therefore requires larger direct systems, additional network
topologies, more model families, and longer stationary trajectories.  Richer
observable state representations could test which hidden variables account
for projected memory, while prospectively chosen alternative coarse grainings
could establish whether the present estimator dependence persists.






\section{Conclusions}
\label{sec:conclusions}

Interacting LLM agents provide a new class of finite many body system in which
microscopic updates are language conditioned decisions rather than fixed
spin-flip rules. Each agent acts from local evidence, delivered neighbor
signals, and bounded private state, while collective organization emerges from
repeated interaction across a network. This preserves the central
statistical mechanical question of how microscopic interactions generate
macroscopic behavior, but places it in a system with contextual,
history-dependent, and partially observed transition rules. The resulting
dynamics can nevertheless be characterized through order, fluctuations,
correlations, response, persistence, and temporal asymmetry. The principal
statistical-mechanical contribution is therefore a reproducible
micro to macro framework for formulating and measuring collective dynamics
when the local transition law is generated by an LLM.

The controlled field quench protocol showed that these measurements reveal
behavior that is not captured by final consensus or population means alone.
Reversing the agents' private evidence produced a supported collective
response across two independently developed model families, while restoring
that evidence revealed model dependent recovery over a fixed observation
horizon. Persistent history also changed the temporal organization of the
projected dynamics relative to both current state only prompting and a
format matched scrambled history control. The latter comparison shows that the
relationship between an agent's recent history and its current decision
matters beyond the generic effect of adding more structured text to a prompt.
Spatial correlations, finite window diagnostics, and the out of sample kinetic
surrogate further identified which aspects of the response admit a reduced
description and which remain dependent on model family, history, or richer
collective coordinates. The weaker Qwen scrambled history decomposition,
Granite's incomplete finite horizon recovery, and the estimator sensitivities
define the boundaries of these supported findings rather than being hidden by
the pooled results.

For agentic AI, these results show why multi-agent systems should be evaluated
as dynamical populations rather than only through individual answer quality,
task completion, or eventual agreement. Systems with similar average outcomes
can differ in internal coordination, sensitivity to changing evidence,
dependence on memory, propagation of local signals, and recovery after
perturbation. Statistical mechanical observables therefore provide a useful
basis for comparing agent architectures, communication rules, memory policies,
and control interventions, while reduced models test whether those collective
dynamics can be predicted from a manageable set of variables. Larger networks,
additional topologies and model families, longer trajectories, and richer
out of sample closures are natural next steps toward determining which
features generalize. The interpretation remains effective rather than
literal: decoding temperature is not physical temperature, reference energy
is not microscopic energy, and projected block-reversal divergence is not
exact thermodynamic entropy production. Within those boundaries, the framework
provides a disciplined way to measure, compare, and ultimately control the
collective behavior of interacting agentic AI systems.






\section*{Data and code availability}

Code, simulation configurations, figure source data, aggregate results, and
reproduction scripts are publicly available at
\url{https://github.com/mantzaris/ThermoAgent}. The reported cross-model
experiment was run on an NVIDIA GeForce RTX 4090 and required 19.19 measured
generation GPU-hours.

Raw prompts, model completions, and complete trajectory records are not
included in the public repository because of their storage volume and the
potential sensitivity of generated content. The original external raw artifact
tree is no longer retained. Content-addressed manifests, replay records, and
aggregate-validation artifacts are provided to document the reconstruction
and verify the reported aggregate results.




\section*{Acknowledgements}

OpenAI Codex (GPT-5) was used for code assistance, the author reviewed all resulting
material and takes responsibility for the manuscript.






\appendix



\section{Study chronology, correction, and inferential details}
\label{app:study-inference}

The earlier quench analysis originally derived each field panel's recovery
threshold from that panel's own baseline, although the protocol called for a
threshold fitted only to the training clusters.  The corrected analysis maps
the leave-one-cluster-out training threshold into the held-out recovery
summary without changing any recorded trajectory.  Figure
\ref{fig:v14recovery} retains every cluster-level recovery path, including the
high-response cluster, under the corrected training-only construction.  The
historical maximum-minus-final recovery statistic remains numerically
reproducible but is not used as directional evidence because its construction
is nonnegative for nearly every nonconstant trajectory.

\resultfigure{figure09_cluster_recovery.pdf}{0.94\textwidth}{%
All field-reversal cluster trajectories from the earlier experiment,
including the high-response cluster.  The corrected audit replaces the
held-out-panel threshold with the training-cluster nominal threshold.  The
historical recovery sign statistic is not used as directional evidence.}{fig:v14recovery}

The delayed audit executes the omitted 10,000 cluster-preserving label
permutations.  All four panels stay together within each cluster, and every
replicate refits imputation, scaling, logistic regression, and leave-one-
cluster-out predictions.  Three-, five-, and seven-sweep nominal geometries
are fitted independently.  Each coordinate and observable family is deleted
in turn, and dependence estimates receive circular-shift null floors.  These
were delayed prespecified sensitivity analyses completed after the primary
outcomes because they had been omitted from the initial implementation, not a
new prospective experiment.

The representation permutation retains the four condition panels in every
cluster and permutes labels only within cluster.  For each of 10,000
replicates, every leave-one-cluster-out fold refits imputation, scaling, and the
multinomial logistic model; test-cluster features do not enter preprocessing.
Figure~\ref{fig:v14audit} records the delayed dependence and full-pipeline
permutation audit.

\resultfigure{figure10_delayed_audit.pdf}{0.94\textwidth}{%
Delayed prespecified dependence and permutation audits.  Raw and
bias-adjusted dependence, window sensitivity, and cluster-preserving nulls are
shown explicitly.  These audits were completed after the primary outcomes
because of implementation omissions.}{fig:v14audit}

For the prospective cross-model tests, the one-sided sign-flip probability
enumerates all $2^n$ sign assignments of the observed cluster effects for
$n\leq20$ paired clusters.  The cluster bootstrap samples complete
graph/environment clusters with replacement and never resamples time points
as independent observations.  H1 receives alpha 0.02; H2--H4 receive a
Holm-adjusted family alpha 0.03.  The allocation was fixed before outcomes
were examined.  Intervals are uncertainty summaries and do not replace the exact
paired tests.





\section{Estimator, control, and closure sensitivities}
\label{app:sensitivities}

Total correlation nulls independently circular shift each coordinate by a
nonzero offset where possible.  This retains marginal occupancy and each
coordinate's internal temporal ordering while breaking synchronous dependence.
The adjusted estimate is not truncated.  Alternative rolling windows rebuild
the words, covariance, training threshold, and held out distance rather than
subsampling a five sweep output.

For block reversal, time shuffled floors quantify finite length occupancy
bias, they do not prove stationarity. Results at block lengths two, three,
and four expose finite history sensitivity, while conditional mutual
information at history depths one through three provides a separate diagnostic
of incomplete first order closure. Figure~\ref{fig:path-sensitivity} displays
the prespecified estimator-sensitivity grid.

\resultfigure{figure11_path_reversal_sensitivity.pdf}{0.94\textwidth}{%
Block-length, pseudocount, and shuffled-bias-floor sensitivity for
bias-adjusted block reversal divergence.  Memory contrasts can be checked
against estimator choices; observable coarse-graining and finite length remain
limiting.}{fig:path-sensitivity}

The scrambled-history arm controls the presence and approximate length of the
history prompt while destroying its temporal relation to the current
trajectory.  Figure~\ref{fig:prompt-balance} records the token-count diagnostic
for that comparison.  It supports approximate prompt-length and format balance
but does not establish turn-by-turn token identity or semantic exchangeability.

\resultfigure{figure12_prompt_balance.pdf}{0.94\textwidth}{%
Scrambled-history prompt-length control.  Per-cluster mean token counts verify
that prompt length and format are approximately matched; semantic content
cannot be exactly token-matched turn by turn.}{fig:prompt-balance}

Finally, CPU trajectories at several sizes provide context for the reduced
kinetic model.  Figure~\ref{fig:surrogate-size} places the direct-system anchor
beside that effective-model size comparison.  These calculations are not
direct-LLM finite-size scaling and do not support a thermodynamic-limit claim.

\resultfigure{figure13_surrogate_size_context.pdf}{0.94\textwidth}{%
Direct $N=16$ anchors in a denser inexpensive effective-model size context.
These are CPU kinetic-surrogate quench responses, not direct-LLM finite-size
scaling.}{fig:surrogate-size}




\section{Finite-window persistence and Binder diagnostics}
\label{app:persistence-binder}

The autocorrelation and Binder diagnostics defined in
equations~\eqref{eq:autocorrelation} and~\eqref{eq:binder} are summarized in
Figure~\ref{fig:persistence-binder}.  Each quantity is computed within one
complete mode graph environment trajectory before uncertainty is summarized
across the six clusters per model.  The autocorrelation sum uses the fixed
two sweep truncation, while one and three sweep truncations remain
machine-readable sensitivity analyses.  Zero-variance phase series are left
undefined rather than assigned artificial persistence, and uncertainty
resamples complete clusters rather than treating time points as independent.

The Binder statistic is likewise calculated trajectory first and describes
the shape of a finite $N=16$ magnetization distribution.  The source data table
retains full phase, early half, and late half estimates under both cluster mean
and pooled moment constructions, near-zero second-moment cases remain
undefined.  Empirical magnetization occupancy is shown alongside the Binder
values so that aggregation sensitive cumulants remain connected to the
finite-window distributions from which they are calculated.

The autocorrelation sum is not an equilibrium correlation time, and these
short, potentially nonstationary windows cannot establish critical slowing
down.  The Binder values are not a crossing analysis and do not establish a
critical point or phase transition.  Figure~\ref{fig:persistence-binder}
therefore presents persistence, Binder, and occupancy together rather than as
an additional confirmatory hypothesis family.

\resultfigure{figure14_persistence_binder.pdf}{0.98\textwidth}{%
Finite-window persistence and order-parameter shape.  Autocorrelation curves
and their 95\% intervals compare the nominal late window with the three field
arms using complete clusters, integrated correlation sums show all cluster
values at the fixed two sweep truncation.  Binder cumulants are paired with
empirical belief magnetization occupancies so their finite sample meaning
remains visible.  These diagnostics do not establish critical slowing down or
a phase transition.}{fig:persistence-binder}




\section{Computational implementation and validation}
\label{app:provenance}

An agent receives only its own persistent state, the delivered neighbor
packets available at its scheduled update, and the condition specific prompt
content defined in Section~\ref{sec:system-protocol}.  Its response changes the
scheduled instance's recorded belief, action, commitment, memory code, signal,
and typed tool consequence, the scheduler validates but does not replace those
fields.  Transition records include the scheduled identity, recipient, packet
bytes, and prompt and output hashes.  Dedicated checks confirm that perturbing
one instance's memory or decision provider leaves unrelated agents' private
state unchanged.

Every experimental arm instantiates a fresh network state from its
prespecified panel seeds.  Model weights and the tokenizer are shared
read-only for throughput, but conversational histories, generation caches,
mutable agent objects, and scientific random-number state are not carried
between arms.  Each generation resets its per-decision seed.  The two pinned
model revisions reported in Section~\ref{sec:cross-model-design} use
Transformers 4.55.4, PyTorch 2.8.0+cu128, bitsandbytes 0.47.0, NF4 double
quantization, and BF16 computation under the common decoding settings given
there.




\bibliographystyle{unsrtnat}
{\small
\bibliography{references}
}




\end{document}
